\section{Consenso y competencia actuales: cómo se ponen al día los equipos chinos}

Luo Fuli (罗福莉) considera que ya existe consenso en que el camino de Anthropic era el correcto; una vez que el camino quedó más claro, los equipos chinos de modelos grandes entraron en una fase de recuperación acelerada. Hoy la brecha generacional en pre-train es muy pequeña, e incluso en algunas arquitecturas los equipos chinos tienen ventaja; donde de verdad hay que ir all in es en el Agent post-train, en particular en cómo hacer bien el RL scaling sobre Agents.

\subsection{Por qué Coding importa en cada paradigma}

Coding ha tocado puntos decisivos en varias etapas. En la etapa de preentrenamiento, los datos de código tienen alta calidad y una estructura clara; en la etapa de reasoning, el código y las matemáticas ofrecen retroalimentación verificable; en la etapa de Agent, el desarrollo de software es en sí mismo una tarea de largo horizonte, con entornos reales, pruebas, depuración y dependencias complejas. Es a la vez señal de entrenamiento y escenario de producto.

\begin{importantbox}{El triple papel de Coding}
Coding es datos, evaluación y producto. Como datos, es estructurado y de alta calidad; como evaluación, cuenta con pruebas y retroalimentación de ejecución; como producto, corresponde directamente a escenarios de productividad de alto valor. Por eso Coding importa en cada paradigma.
\end{importantbox}

\subsection{El entorno importa más que la experiencia}

En la parte final, Luo Fuli enfatiza que el entorno importa más que la experiencia. Por entorno entiende la capacidad de cómputo, la infraestructura, los chips, el sistema operativo, el tráfico, las redes sociales, la cultura organizacional y los recursos estratégicos. Que un equipo pueda liderar en el nuevo paradigma depende de si logra que todos estos recursos se adapten en conjunto a la arquitectura de Agent.

\begin{warningbox}{La experiencia puede convertirse en una carga}
Cuando el paradigma cambia con rapidez, la experiencia previa ayuda a juzgar, pero también limita la imaginación. Lo que de verdad importa es si el entorno le da al equipo retroalimentación y aceleración continuas, de modo que pueda refutar una y otra vez sus juicios anteriores y reconstruir sus procesos anteriores.
\end{warningbox}

\subsection{Resumen de la sección}

El consenso de la competencia actual es: la brecha en pre-train se reduce y el Agent post-train se vuelve la línea principal. La oportunidad de los equipos chinos está en que el camino es más claro, son sensibles al costo y la retroalimentación de las aplicaciones es rápida; los retos están en la RL infra, los entornos de tareas y la reestructuración organizacional.
